% Shared preamble for the supporting-information documents (S1-S5).
\documentclass[10pt,letterpaper]{article}
\usepackage[margin=1in]{geometry}
\usepackage{amsmath,amssymb}
\usepackage{booktabs,array,longtable}
\usepackage[table]{xcolor}
\usepackage{hyperref}
\usepackage{microtype}
\usepackage{listings}
\usepackage{graphicx}
\usepackage[labelfont=bf,labelsep=period,justification=raggedright,singlelinecheck=off]{caption}
\definecolor{linkblue}{HTML}{1C5CAB}
\definecolor{codebg}{HTML}{F4F4F1}
\hypersetup{colorlinks=true, linkcolor=linkblue, citecolor=linkblue, urlcolor=linkblue}
\lstset{basicstyle=\ttfamily\footnotesize, backgroundcolor=\color{codebg},
  frame=none, breaklines=true, columns=fullflexible, keepspaces=true,
  xleftmargin=4pt, xrightmargin=4pt, aboveskip=6pt, belowskip=6pt}
\setlength{\parindent}{0pt}
\setlength{\parskip}{5pt}
\raggedright
\renewcommand{\tablename}{Table}
\newcommand{\sititle}[2]{{\Large\bfseries #1}\par\vspace{3pt}{\large #2}\par\vspace{2pt}
  {\small\color{black!60}Supporting information for ``Environmental fluctuation limits gene regulation unless it can be tracked''}\par\vspace{10pt}}

\begin{document}
\sititle{S5 Table.}{Prior work and how this study relates to it}

The table sets out which parts of the argument are established in the
literature, which conflict with our results, and which we did not find
precedent for. Literature searches covered theory of plasticity and cue
reliability, gene-regulatory-network evolution under fluctuating selection,
digital evolution (Avida, Aevol), network sparseness and modularity, and weighted
network measures.

\begin{longtable}{@{}p{0.3\linewidth}p{0.18\linewidth}p{0.46\linewidth}@{}}
\toprule
Claim & Status & Sources and relation to this study\\
\midrule
\endhead
Reliable cues favour plastic responses; without them, fluctuation favours
generalists, bet-hedging or genetic tracking & Established &
Levins 1968 [1]; Moran 1992 [2]; Kussell and Leibler 2005 [3]; Botero et al.
2015 [4]. Our results concern the corresponding effect on the regulatory network.\\
In a gene regulatory network model, plasticity evolves mainly under fast, erratic
change with reliable cues & Established & Jorritsma and van den Berg 2026 [5].
Their abstract reports which parts of the network become sensitive to mutation.
We have not been able to read the full text; if it reports connection strength
by cue reliability, part of our architectural result would overlap with it.\\
Fluctuating environments can select for or against gene regulation &
Established experimentally & Pollack, Nozoe and Kussell 2025 [6], who report that catabolite symport flux strongly affects selection for regulation and note that laboratory evolution under metabolic fluctuation often loses the regulation of metabolic genes.\\
Fluctuation without information about the optimum promotes complex regulatory
networks & Published; conflicts with our R-Aevol results for a random-walk
optimum & Tsuda and Kawata 2010 [7]. Their optimum moved across generations, and
the networks received an on/off signal that carried no information about its
position. The effect was strong for a random walk and slight, limited to large
and infrequent shifts, for an optimum alternating between two positions, which
is the pattern of our design. Our grain experiment found that grain does not
reconcile the two results. The remaining candidate explanations are the novel
optima of the random walk and differences in the genotype--phenotype map.\\
Fluctuating environments favour larger genomes & Published; matches earlier EvoSim results & Bentkowski et al. 2015 [8]. R-Aevol shows the opposite.\\
Genome size is limited by selection for mutational robustness and falls with
population size and mutation rate & Established in Aevol & Knibbe et al. 2007 [9];
Luiselli et al. 2024 [10]. Consistent with the small genomes in large EvoSim
populations.\\
Sensing and non-sensing populations evolve differently in fluctuating
environments & Published, different measures & Lalejini et al. 2021 [11]
(Avida), measuring trait loss and the accumulation of deleterious mutations. Network architecture was not measured.\\
Regulatory complexity is measured by counting genes or edges & Common practice &
Tsuda and Kawata 2010 [7] (gene numbers, degree distributions); Espinosa-Soto
2018 [12] (binary edges, weights not discussed).\\
Network size and connectivity can grow without selection for them & Established &
Soyer and Bonhoeffer 2006 [13].\\
The number of transcription factors scales faster than genome size and is
higher in free-living bacteria & Established; two interpretations &
Comparative genomics links it to lifestyle and environmental variety (van
Nimwegen 2003 [16]; Cases, de Lorenzo and Ouzounis 2003 [17]). R-Aevol
simulations reproduce the scaling and attribute it to mutation and
rearrangement rates instead of lifestyle (Beslon et al. 2010 [18]).\\
Summed edge strength as a network measure & Standard method & Iwasaki, Tsuda and
Kawata 2013 [14].\\
The yoked control & Established in the study of learning & Church 1964 [15]
describes a bias that arises when individual subjects are paired; our arms are
not paired, so it does not apply.\\
\midrule
\multicolumn{3}{@{}l}{\emph{Not found in the literature}}\\
Edge count does not track edge strength under fluctuation & This study &
EvoSim (exploratory; more edges with the same total strength), R-Aevol
informative-cue populations (pre-registered secondary analysis; no more links
but stronger regulation).\\
A yoked control that separates a cue's information from its presence in the
evolution of regulatory networks & This study & R-Aevol, with switching rates
matched between arms (S2 Text).\\
Regulatory strength and genome size compared across informative, yoked and
absent cues & This study & Pre-registered predictions P1--P4; subject to the
caveat on [5] above.\\
\bottomrule
\end{longtable}

\section*{References}
\small
\begin{enumerate}\setlength{\itemsep}{1pt}
\item Levins R. Evolution in Changing Environments. Princeton: Princeton University Press; 1968.
\item Moran NA. The evolutionary maintenance of alternative phenotypes. Am Nat. 1992;139(5):971--89.
\item Kussell E, Leibler S. Phenotypic diversity, population growth, and information in fluctuating environments. Science. 2005;309(5743):2075--8.
\item Botero CA, Weissing FJ, Wright J, Rubenstein DR. Evolutionary tipping points in the capacity to adapt to environmental change. Proc Natl Acad Sci USA. 2015;112(1):184--9.
\item Jorritsma A, van den Berg P. The evolution of plasticity and evolvability in a simple gene regulatory network. Evolution. 2026;80(1):230--9.
\item Pollack D, Nozoe T, Kussell E. Gene regulation, phenotypic memory, and selection in fluctuating environments. bioRxiv [Preprint]. 2025. doi:10.1101/2025.06.23.661119.
\item Tsuda ME, Kawata M. Evolution of gene regulatory networks by fluctuating selection and intrinsic constraints. PLoS Comput Biol. 2010;6(8):e1000873.
\item Bentkowski P, van Oosterhout C, Mock T. A model of genome size evolution for prokaryotes in stable and fluctuating environments. Genome Biol Evol. 2015;7(8):2344--51.
\item Knibbe C, Coulon A, Mazet O, Fayard JM, Beslon G. A long-term evolutionary pressure on the amount of noncoding DNA. Mol Biol Evol. 2007;24(10):2344--53.
\item Luiselli J, Rouzaud-Cornabas J, Lartillot N, Beslon G. Genome streamlining: effect of mutation rate and population size on genome size reduction. Genome Biol Evol. 2024;16(12):evae250.
\item Lalejini A, Ferguson AJ, Grant NA, Ofria C. Adaptive phenotypic plasticity stabilizes evolution in fluctuating environments. Front Ecol Evol. 2021;9:715381.
\item Espinosa-Soto C. On the role of sparseness in the evolution of modularity in gene regulatory networks. PLoS Comput Biol. 2018;14(5):e1006172.
\item Soyer OS, Bonhoeffer S. Evolution of complexity in signaling pathways. Proc Natl Acad Sci USA. 2006;103(44):16337--42.
\item Iwasaki WM, Tsuda ME, Kawata M. Genetic and environmental factors affecting cryptic variations in gene regulatory networks. BMC Evol Biol. 2013;13:91.
\item Church RM. Systematic effect of random error in the yoked control design. Psychol Bull. 1964;62:122--31.
\item van Nimwegen E. Scaling laws in the functional content of genomes. Trends Genet. 2003;19(9):479--84.
\item Cases I, de Lorenzo V, Ouzounis CA. Transcription regulation and environmental adaptation in bacteria. Trends Microbiol. 2003;11(6):248--53.
\item Beslon G, Parsons DP, Sanchez-Dehesa Y, Pe\~na JM, Knibbe C. Scaling laws in bacterial genomes: a side-effect of selection of mutational robustness? BioSystems. 2010;102(1):32--40.
\end{enumerate}

\end{document}
